\section{El teorema de las cadenas pequeñas}

% ------------------------------------------------------------------------
\subsection{Enunciado}

\begin{frame}
  \frametitle{El teorema de las cadenas pequeñas}

  \begin{theorem}
    Sea $\cU$ un recubrimiento de un espacio $X$ e
    $\iota\colon\CU_\bullet(X)\hookrightarrow\Sing_\bullet(X)$ la inclusión.
    Existen un morfismo de complejos
    $\rho\colon\Sing_\bullet(X)\to\CU_\bullet(X)$ y homomorfismos
    $D\colon\Sing_n(X)\to\Sing_{n+1}(X)$ tales que
    \begin{enumerate}
      \item $\rho\iota=\id$, es decir, $\rho$ fija $\CU_\bullet(X)$;
      \item $\partial D+D\partial=\id-\iota\rho$;
      \item $D$ se anula sobre $\CU_\bullet(X)$.
    \end{enumerate}
  \end{theorem}

  En particular, $\iota$ es una equivalencia de homotopía de cadenas con
  inversa homotópica~$\rho$.

  \medskip

  La demostración se hará en varias etapas, usando la subdivisión iterada y el lema del número de Lebesgue.
\end{frame}


% ------------------------------------------------------------------------
\subsection{El número de subdivisiones}

\begin{frame}[allowframebreaks]
  \frametitle{El número de subdivisiones de un símplice}

  Sea $\sigma\colon\simplex{n}\to X$ un símplice singular. Por la
  proposición de subdivisión suficiente, existe un menor entero $k(\sigma)\geq0$ tal que
  $\Sd^i\sigma\in\CU_n(X)$ para todo $i\geq k(\sigma)$.

  \begin{definition}
    Definimos $m(\sigma)\geq0$ por recurrencia en $n$:
    \[
      m(\sigma)=
      \begin{cases}
        k(\sigma), & n=0,\\
        \max\bigl\{k(\sigma),m(\partial_0\sigma),\ldots,m(\partial_n\sigma)\bigr\},
          & n\geq1.
      \end{cases}
    \]
  \end{definition}

  \framebreak

  \begin{lemma}
    Para todo símplice singular $\sigma\colon\simplex{n}\to X$:
    \begin{enumerate}
      \item $\Sd^i\sigma\in\CU_n(X)$ para todo $i\geq m(\sigma)$;
      \item si $n\geq1$, $m(\partial_j\sigma)\leq m(\sigma)$ para todo $j$;
      \item si $\sigma\in\CU_n(X)$, entonces $m(\sigma)=0$.
    \end{enumerate}
  \end{lemma}

  \begin{proof}
    Las dos primeras propiedades se siguen de la definición, ya que
    $m(\sigma)\geq k(\sigma)$ y $m(\sigma)\geq m(\partial_j\sigma)$.

    \medskip

    La tercera, por inducción en $n$. Como $\Sd$ preserva
    $\CU_\bullet(X)$, $\Sd^i\sigma\in\CU_n(X)$ para todo $i\geq0$, luego
    $k(\sigma)=0$. Si $n\geq1$, las caras $\partial_j\sigma$ también están en
    $\CU_{n-1}(X)$, y por hipótesis de inducción $m(\partial_j\sigma)=0$.
    Así, $m(\sigma)=0$.
  \end{proof}
\end{frame}

% ------------------------------------------------------------------------
\subsection{Construcción de la retracción}

\begin{frame}[allowframebreaks]
  \frametitle{Definición de la homotopía y de la retracción}

  \begin{definition}
    Sea $D\colon\Sing_n(X)\to\Sing_{n+1}(X)$ el homomorfismo dado sobre los
    símplices singulares por
    \[
      D\sigma=T_{m(\sigma)}\sigma,
    \]
    y sea \[\rho=\id-\partial D-D\partial\colon\Sing_n(X)\to\Sing_n(X).\]
  \end{definition}

  \framebreak

  \begin{lemma}
    El homomorfismo $\rho\colon\Sing_\bullet(X)\to\Sing_\bullet(X)$ es un morfismo de complejos, $\partial\rho=\rho\partial$.
  \end{lemma}

  \begin{proof}
    Como $\partial^2=0$,
    \[
      \partial\rho
      = \partial(\id-\partial D-D\partial)
      =\partial-\partial D\partial
      \qquad\text{y}\qquad
      \rho\partial
      =(\id-\partial D-D\partial)\partial
      =\partial-\partial D\partial.
    \]
  \end{proof}

  Por definición, $\partial D+D\partial=\id-\rho$. Falta ver que $\rho$
  toma valores en $\CU_\bullet(X)$ y las demás propiedades.
\end{frame}

\begin{frame}
  \frametitle{La homotopía se anula sobre las cadenas subordinadas}

  \begin{lemma}
    El homomorfismo $D$ se anula sobre $\CU_\bullet(X)$ y $\rho$ fija
    $\CU_\bullet(X)$.
  \end{lemma}

  \begin{proof}
    Sea $\sigma\colon\simplex{n}\to X$ un símplice singular subordinado a $\cU$. Por un lema anterior,
    $m(\sigma)=0$, luego $D\sigma=T_0\sigma=0$. Esto prueba que $D$ se anula
    sobre $\CU_\bullet(X)$.

    \medskip

    Además, $\partial\sigma\in\CU_{n-1}(X)$, así que también
    $D\partial\sigma=0$, y por tanto
    \[
      \rho\sigma=\sigma-\partial D\sigma-D\partial\sigma=\sigma.
    \]
  \end{proof}
\end{frame}

\begin{frame}[allowframebreaks]
  \frametitle{La retracción toma valores en las cadenas subordinadas}

  \begin{lemma}
    Se tiene que $\rho\bigl(\Sing_n(X)\bigr)\subseteq\CU_n(X)$.
  \end{lemma}

  \begin{proof}
    Sea $\sigma\colon\simplex{n}\to X$ un símplice singular. Si $n=0$, la
    imagen de $\sigma$ es un punto, que está en algún $U_i$, luego
    $\sigma\in\CU_0(X)$ y $\rho\sigma=\sigma$ por el lema anterior.

    \medskip

    Si $n\geq1$, sean $m=m(\sigma)$ y $m_j=m(\partial_j\sigma)$, que cumplen
    $m_j\leq m$. Por el lema de iteración de homotopías,
    \[
      \partial D\sigma=\partial T_m\sigma
      =\sigma-\Sd^m\sigma-T_m\partial\sigma.
    \]

    \framebreak

    Sustituyendo en la definición de $\rho$ y usando que
    $\partial\sigma=\sum_j(-1)^j\partial_j\sigma$,
    \begin{align*}
      \rho\sigma
      &=\sigma-\partial D\sigma-D\partial\sigma
       =\Sd^m\sigma+T_m\partial\sigma-D\partial\sigma\\
      &=\Sd^m\sigma+\sum_{j=0}^n(-1)^j\bigl(T_m-T_{m_j}\bigr)(\partial_j\sigma)\\
      &=\Sd^m\sigma+\sum_{j=0}^n(-1)^j\sum_{i=m_j}^{m-1}T\Sd^i(\partial_j\sigma).
    \end{align*}

    El primer término está en $\CU_n(X)$ porque $m=m(\sigma)$. En los
    demás, $i\geq m_j=m(\partial_j\sigma)$, luego
    $\Sd^i(\partial_j\sigma)\in\CU_{n-1}(X)$, y como $T$ preserva las cadenas
    subordinadas, $T\Sd^i(\partial_j\sigma)\in\CU_n(X)$. Por tanto
    $\rho\sigma\in\CU_n(X)$.
  \end{proof}
\end{frame}

% ------------------------------------------------------------------------
\subsection{Consecuencias}

\begin{frame}
  \frametitle{Demostración del teorema}

  Los lemas anteriores prueban el teorema:
  \begin{itemize}
    \item $\rho\colon \Sing_\bullet(X)\to\CU_\bullet(X)$ es un morfismo de complejos;
    \item $\rho$ fija $\CU_\bullet(X)$, es decir, $\rho\iota=\id$;
    \item $\partial D+D\partial=\id-\iota\rho$, por definición de $\rho$;
    \item $D$ se anula sobre $\CU_\bullet(X)$.
  \end{itemize}
\end{frame}

\begin{frame}
  \frametitle{Consecuencia en homología}

  \begin{corollary}
    La inclusión $\iota\colon\CU_\bullet(X)\hookrightarrow\Sing_\bullet(X)$
    induce isomorfismos
    \[
      H_n\bigl(\CU_\bullet(X)\bigr)\cong H_n(X)
    \]
    para todo $n$, con inverso inducido por $\rho$.
  \end{corollary}

  \medskip

  Es decir, la homología singular puede calcularse usando solo símplices
  subordinados a un recubrimiento cualquiera.
\end{frame}

\begin{frame}[allowframebreaks]
  \frametitle{Versión reducida}

  El complejo aumentado $\tCU_\bullet(X)$ coincide con $\CU_\bullet(X)$
  en grados $n\geq0$ y con $\widetilde{\Sing}_{-1}(X)=\Z\,[\;]$ en grado
  $-1$. Es un subcomplejo de $\widetilde{\Sing}_\bullet(X)$.

  \begin{corollary}
    La inclusión $\tCU_\bullet(X)\hookrightarrow\widetilde{\Sing}_\bullet(X)$
    induce isomorfismos
    \[
      H_n\bigl(\tCU_\bullet(X)\bigr)\cong\widetilde{H}_n(X)
    \]
    para todo $n$.
  \end{corollary}

  \framebreak

  \begin{proof}
    Extendemos $\rho$ por la identidad y $D$ por $0$ en grado $-1$. Todo
    $0$-símplice tiene imagen en un punto, que está en algún $U_i$, así que
    $\CU_0(X)=\Sing_0(X)$ y $\rho$ es la identidad en grado $0$. Por tanto
    $\rho$ conmuta con el borde $\Sing_0(X)\to\widetilde{\Sing}_{-1}(X)$, y
    es un morfismo de complejos aumentados.

    \medskip

    Las identidades $\rho\iota=\id$ y $\partial D+D\partial=\id-\iota\rho$
    se cumplen en grados $n\geq1$ como antes. En grado $0$ también, pues
    $D\partial\sigma=D[\;]=0$, y en grado $-1$ ambos miembros son nulos.
    Así, la inclusión es una equivalencia de homotopía de cadenas.
  \end{proof}
\end{frame}
